\scanpage{132}
\mspage{-III- \handwritten{4} -}
\begin{flushright}number \handwritten{383}\end{flushright}
$s_i$, hence for $G$ (proposition 1), and is consequently orthogonal to $W$ ; this shows that the $e_i$ span $W$. Which was to be proved.

This being so, let $I=\handwritten{\{}1,2,\ldots,m\handwritten{\}}$, $M=(m_{ij})$ be the Coxeter matrix on $I$ where $m_{ij}$ is the order of $s_i s_j$ in $G$, and $G(M)$ the Coxeter group defined by the generators $r_i$ $(i\in I)$ and the relations $(r_i r_j)^{m_{ij}}=1$ (\S\ 2, number 1). By proposition 3 one can take $W$, equipped with the basis $(e_i)$, as the space of the direct representation $\handwritten{\rho}$ of $G(M)$ ; by proposition 2 and the fact that the $e_i$ are unit vectors, the bilinear form $B_M$ of \S\ 2, number 2 is the restriction to $W$ of the given form $(x\handwritten{|}y)$ ; thus $\handwritten{\rho}(r_i)$ is the restriction $s'_i$ to $W$ of the reflection $s_i$. Let us identify $W$ with its dual by means of the form $(x\handwritten{|}y)$ ; since, to pass from a symmetry to its contragredient, one interchanges vector and linear form (\S\ 2, number 3), and since $s'_i$ is defined by a unit vector and the orthogonal hyperplane, $s'_i$ is its own contragredient. If $(e'_i)$ is the basis of $W$ dual to $(e_i)$ (i.e. $(e_i\handwritten{|}e'_j)=\handwritten{\delta}_{ij}$ for all $i,j$), the set of $\displaystyle x=\handwritten{\sum}_{i\in I}a_i e'_i$ $(a_i\in R)$ in $W$ such that $(x\handwritten{|}e_i)>0$ for all $i$ is that of the $x$ such that $a_i>0$ for all $i$ ; hence, by (2), $C\handwritten{\cap}W$ is the chamber denoted by $C$ in \S\ 2, number 3. Taking into account the fact that $G$ acts trivially on the subspace $W'$ orthogonal to $W$, theorem 1 of \S\ 2, number 3 gives the following result :

THEOREM 1.- \uline{Let} $\underline{H}$ \uline{be a set of hyperplanes of} $E$ \uline{satisfying} (1) \uline{and} (2), $C$ \uline{a chamber defined by} $\underline{H}$, $H_1,\ldots,H_m$ \uline{its walls}, \uline{and} $s_i$ \uline{the reflection with respect to} $H_i$. \uline{Write} $m_{ij}$ \uline{for the order of} $s_i s_j$ \uline{in} $G$. \uline{Then}

a) \uline{the hyperplanes} $H_i$ \uline{are linearly independent} ;

b) \uline{the group} $G$ \uline{is defined by the generators} $s_i$ \uline{and the relations} $s_i^2=(s_i s_j)^{m_{ij}}=1$;

c) \uline{the group} $G$ \uline{acts simply transitively on the set of chambers defined by} $\underline{H}$.

COROLLARY 1.- \uline{If} $G$ \uline{is essential}, \uline{any two vectors} $x,y$ \uline{in the same chamber} $C$ \uline{form an acute angle} (i.e. $(x\handwritten{|}y)\handwritten{\geq}0$).
% archive sheet number: 131
